\paragraph{}Este capítulo introduz o leitor ao contexto no qual a presente pesquisa se insere, apresentando as motivações iniciais, o escopo do problema de otimização selecionado e a relevância científica da abordagem proposta. A \autoref{sec:tema} delineia o tema central do estudo, abordando as restrições e a hierarquia do Problema de Roteamento de Veículos com Janelas de Tempo (VRPTW). Em seguida, a \autoref{sec:delimitacao} estabelece as delimitações do trabalho, focando no desenvolvimento e avaliação empírica de meta-heurísticas integradas em um sistema multi-agente colaborativo guiado por aprendizado por reforço. A justificativa para a aplicação destas ferramentas algorítmicas, dadas as dificuldades inerentes aos métodos exatos diante de problemas classificados como \textit{NP-hard}, é apresentada na \autoref{sec:justificativa}. Por fim, as Seções \ref{sec:objetivos}, \ref{sec:metodologia} e \ref{sec:descricao} expõem, respectivamente, os objetivos formulados, a metodologia adotada de melhoria contínua e a estruturação capitular do conteúdo subsequente neste documento.

\section{Tema}
\label{sec:tema}

\paragraph{}O tema do trabalho é a aplicação e a análise de métodos meta-heurísticos para a resolução de problemas de otimização complexos, classificados como \textit{NP-hard} \cite{Garey79}, como o Problema de Roteamento de Veículos (ou \textit{Vehicle Routing Problem} -- VRP). O VRP, introduzido pioneiramente na literatura sob a ótica da distribuição de derivados de petróleo em 1959 \cite{Dantzig59}, consiste em determinar o conjunto ideal de rotas para uma frota de veículos, com o objetivo de visitar um conjunto de clientes, cada um com uma demanda específica. O problema a ser resolvido é encontrar a solução mais eficiente para este roteamento, considerando restrições como distância total, número de veículos, tempo e capacidade de entrega.

\paragraph{}Dentro do vasto domínio do VRP, este trabalho concentra-se especificamente no Problema de Roteamento de Veículos com Janelas de Tempo (\textit{Vehicle Routing Problem with Time Windows} -- VRPTW). O VRPTW estende o problema clássico ao impor que o atendimento a cada cliente deve ser iniciado obrigatoriamente dentro de um intervalo temporal pré-definido, composto por um tempo de abertura e um tempo de fechamento (\textit{time window}).



\paragraph{}A inclusão destas restrições temporais torna o desafio ainda mais rigoroso, pois a viabilidade de uma rota não depende apenas da capacidade de carga do veículo, mas da coordenação precisa dos horários de chegada e dos tempos de serviço em cada nó. Neste cenário, o objetivo de encontrar a solução mais eficiente desdobra-se em uma hierarquia de otimização: busca-se, primordialmente, a minimização do número de veículos necessários e, subsequentemente, a redução da distância total percorrida, garantindo que nenhum compromisso temporal seja violado. Por ser um problema de natureza \textit{NP-hard}, a exploração de janelas de tempo exige que os métodos meta-heurísticos aplicados possuam alta capacidade de adaptação para navegar em um espaço de busca repleto de restrições de inviabilidade.

\section{Delimitação}
\label{sec:delimitacao}

\paragraph{}O objeto de estudo são problemas de otimização \textit{NP-hard}. O projeto se concentrará na otimização do VRPTW, que estende o clássico Problema do Caixeiro Viajante (\textit{Traveling Salesman Problem} -- TSP) para múltiplos veículos, capacidade de carga e janelas de tempo. O foco será na implementação e comparação de diferentes algoritmos meta-heurísticos (\textit{Simulated Annealing} -- SA, \textit{Tabu Search} -- TS e \textit{Genetic Algorithm} -- GA), e na integração desses métodos com um sistema multiagente (\textit{Multi-Agent System} -- MAS), conjuntamente com a aplicação de aprendizado por reforço (\textit{Reinforcement Learning} -- RL) para melhoria deste sistema.

\section{Justificativa}
\label{sec:justificativa}

\paragraph{}O VRPTW, em virtude de sua formulação que estende o VRP clássico, é rigorosamente classificado como um problema \textit{NP-hard} \cite{Garey79}. Nesse contexto de alta complexidade computacional, o esforço temporal demandado por métodos exatos para a obtenção de uma solução ótima cresce de forma exponencial em função da dimensionalidade da instância analisada, tornando tais abordagens intratáveis para a ampla escala de cenários reais. Diante dessa limitação, as meta-heurísticas despontam como uma alternativa viável, visto que, embora não assegurem a otimalidade global, oferecem mecanismos para explorar o espaço de busca eficientemente, fornecendo aproximações de alta qualidade em tempo computacional factível.

\paragraph{}A originalidade deste trabalho reside na aplicação de um MAS com RL para integrar e aprimorar o desempenho dos diferentes métodos meta-heurísticos. O objetivo é explorar como a comunicação entre agentes pode levar a resultados superiores aos de cada algoritmo individualmente, evitando mínimos locais e alcançando uma solução globalmente mais eficiente. A importância do projeto está em sua aplicabilidade prática em áreas como logística, transporte e serviços de entrega, onde a otimização de rotas é fundamental para a eficiência e redução de custos.

\paragraph{}Ainda no tocante à justificativa da abordagem, o Apêndice B apresenta uma análise quantitativa da grandeza temporal requerida pela força bruta para a obtenção da solução ótima nas instâncias discutidas, reiterando a necessidade de métodos heurísticos.

\section{Objetivos}
\label{sec:objetivos}

\paragraph{}O objetivo geral é analisar e implementar um sistema híbrido, combinando métodos meta-heurísticos com um sistema multiagente e aprendizado por reforço, para otimizar a resolução do VRPTW. Como objetivos específicos, temos:

\begin{enumerate}
    \item Formular o VRPTW, definindo a função de custo e as restrições (tempo, volume, peso) para a avaliação das soluções;
    \item Implementar, de forma independente, os algoritmos meta-heurísticos: SA, TS e GA;
    \item Integrar os algoritmos em um MAS, no qual a interação entre os agentes leva à colaboração na busca pela solução;
    \item Aplicar RL para aprimorar a estratégia de busca dos agentes;
    \item Comparar o desempenho dos algoritmos individuais com o MAS e o sistema com aprendizado por reforço (MAS-RL) para avaliar a melhoria na otimização da solução.
\end{enumerate}

\section{Metodologia}
\label{sec:metodologia}

\paragraph{}A metodologia será dividida nas seguintes etapas, focando na aplicação de uma abordagem iterativa e de melhoria contínua para a otimização do problema.

\begin{itemize}
    \item \textbf{Modelagem do Problema:} O VRP será formulado como um problema de otimização, no qual uma solução é representada como um vetor ordenado de pontos de entrega. A função de custo avaliará a qualidade de cada solução, considerando a distância percorrida, o número de veículos e as restrições de tempo e capacidade, atribuindo pesos a cada parâmetro;
    \item \textbf{Implementação dos Algoritmos:} Os três algoritmos meta-heurísticos serão implementados em \textit{Python}. Para cada algoritmo, serão definidas etapas de inicialização de uma solução aleatória, geração de novas soluções e a manutenção de um histórico para direcionar a busca;
    \item \textbf{Integração com Sistema Multiagente:} Os algoritmos serão encapsulados como agentes dentro de um MAS. O sistema operará de forma cooperativa, permitindo que os agentes interajam para compartilhar informações e refinar as soluções encontradas;
    \item \textbf{Aplicação de Aprendizado por Reforço:} O método de \textit{Q-Learning} será aplicado para associar os estados do problema às ações (ajuste de estratégias de busca), permitindo que o sistema aprenda a selecionar a melhor abordagem de forma autônoma ao longo do tempo.
\end{itemize}



\section{Descrição}
\label{sec:descricao}

\paragraph{}O restante deste trabalho está organizado em quatro capítulos principais, além das referências bibliográficas e apêndices, conforme detalhado a seguir.

\paragraph{}No Capítulo 2, é apresentada a Fundamentação Teórica, que estabelece a base conceitual necessária para a compreensão do projeto. Nele, define-se o VRPTW e detalha-se a metodologia de avaliação baseada no \textit{Benchmark} de Solomon \cite{Solomon87}. O capítulo descreve o funcionamento técnico do trio de meta-heurísticas selecionado e introduz a arquitetura do MAS utilizando o \textit{framework} Mesa. Finaliza abordando as estratégias de cooperação via \textit{Path Relinking} (PR) e a teoria de Aprendizado por Reforço.

\paragraph{}O Capítulo 3 detalha o desenvolvimento do Método Proposto da pesquisa. Ele descreve as etapas iniciais de modelagem e as arquiteturas das classes. Em seguida, expõe-se a modelagem do sistema multiagente, demonstrando a dinâmica de cooperação por meio do \textit{Path Relinking} e a implementação do \textit{Q-Learning} como hiper-heurística, orquestrando e selecionando as melhores ações para esse ambiente.

\paragraph{}Os testes realizados e a apresentação dos resultados ficam a encargo do Capítulo 4: Experimentos Computacionais. A seção avalia o desempenho prático, apresentando inicialmente o \textit{baseline} de cada meta-heurística isolada, em sequência os resultados agregados pela cooperação simplista do sistema multiagente e, por fim, a análise do desempenho da hiper-heurística com Aprendizado por Reforço (MAS-RL).

\paragraph{}As reflexões e fechamentos a respeito da pesquisa são apresentados no Capítulo 5, que compõe as Conclusões do trabalho. Nele, será explicitado se a integração entre MAS e RL de fato superou o desempenho dos algoritmos isolados, validando a hipótese de que a inteligência coletiva e a adaptação dinâmica de parâmetros produzam soluções superiores para o VRPTW. O capítulo encerra com a indicação de oportunidades de melhoria e sugestões para trabalhos futuros.